\documentclass[11pt,UTF8,fontset=fandol]{ctexart}
\usepackage{amsmath,amssymb,amsfonts}
\usepackage{graphicx,float,caption,subcaption}
\usepackage{booktabs,array,multirow,makecell}
\usepackage{geometry,enumitem,xcolor}
\usepackage[hidelinks]{hyperref}
\usepackage{url}
\geometry{margin=2.3cm}
\graphicspath{{/Users/yye/project/论文/assets/graphwiz/}}
\captionsetup{font=small,labelfont=bf}
\setlength{\parskip}{0.4em}
\hypersetup{pdftitle={GraphWiz 中文精读导读},pdfauthor={Chen et al.}}

% 小白知识卡片：上下加线 + 粗体标签，稳健不挑内容
\newenvironment{tip}[1]{%
  \par\smallskip\noindent\rule{\linewidth}{0.4pt}\par\nobreak
  \noindent\textbf{【小白补丁】#1}\par\nobreak\noindent\ignorespaces
}{%
  \par\noindent\rule{\linewidth}{0.4pt}\par\smallskip
}

\title{GraphWiz 中文精读导读：教大模型“一步一步”解图论题的开源专家}
\author{原文：Nuo Chen, Yuhan Li, Jianheng Tang, Jia Li（KDD 2024）}
\date{}

\begin{document}
\maketitle
\begin{center}\small\textit{本导读为面向初学者的中文学习笔记：含要点翻译、公式精讲与图表解读。}\end{center}

\section*{论文卡片}
\begin{center}
\begin{tabular}{@{}ll@{}}
\toprule
\textbf{标题} & GraphWiz: An Instruction-Following Language Model for \\
            & Graph Computational Problems \\
\textbf{作者} & Nuo Chen, Yuhan Li, Jianheng Tang, Jia Li（陈诺、李宇晗、唐建衡、李佳） \\
\textbf{机构} & 香港科技大学（广州）数据科学与分析学域 / 香港科技大学 \\
\textbf{会议} & KDD 2024（ACM SIGKDD 知识发现与数据挖掘会议） \\
\textbf{arXiv} & 2402.16029 \\
\textbf{项目页} & \url{https://graph-wiz.github.io} \\
\textbf{代码} & \url{https://github.com/nuochenpku/Graph-Reasoning-LLM} \\
\textbf{方向} & 图（Graph）$\times$ 大模型 $\times$ 指令微调 / 偏好对齐（DPO） \\
\textbf{一句话} & 先造一个带“显式推理过程”的大规模图论指令数据集 GraphInstruct，\\
              & 再用它微调出能解 9 类图论题、且超过 GPT-4 的开源模型 GraphWiz。 \\
\bottomrule
\end{tabular}
\end{center}

\section{一、这篇论文在解决什么问题}

大语言模型（LLM，如 GPT-4、LLaMA）在文字处理上已经很强，但作者发现：让它们真正\textbf{读懂一张图、
并算出图论问题的答案}，这件事还很少有人系统研究。这里说的“图”是数学意义上的图——由\textbf{节点（node）}
和\textbf{边（edge）}组成的结构，比如社交网络、电路、地图路网。而“图计算问题（graph computational
problem）”指的是要靠图算法才能解决的题目，例如：两个节点之间有没有路径？最短路是多少？图里有没有环？

\textbf{已有做法的三个痛点}（作者在引言里点名）：
\begin{enumerate}[leftmargin=2em]
  \item \textbf{只会“猜”，不会“讲”}：以往多用精心设计的提示词（prompt）去“哄” GPT-4 这类闭源模型答题，
        但模型只给一个 yes/no，看不到推理过程——你无法判断它是\emph{真的推理出来}，还是\emph{蒙对了}。
  \item \textbf{开源模型很弱}：同样的提示法，开源的 LLaMA 表现远不如 GPT-4；而且把每条边写成一句话会让
        输入变得很长，限制了能处理的图规模（以往研究普遍把图限制在 20 个节点、100 条边以内）。
  \item \textbf{缺一套“训练教材”}：市面上没有一个开源数据集，能教模型\textbf{把图论题的推理过程一步步写出来}。
\end{enumerate}

\textbf{本文目标}：训练\textbf{单独一个}开源模型，让它既能解多种类型、多种规模的图论问题，又能像图~\ref{fig:fig1}
那样\textbf{用自然语言把推理路径明明白白写出来}。为此作者做了三件事：①造数据集 GraphInstruct；
②用它微调出 GraphWiz；③再用 DPO 进一步对齐，得到更强的 GraphWiz-DPO。

\begin{figure}[H]\centering
  \includegraphics[width=0.92\linewidth]{cand_p02_k01.pdf}
  \caption{原文 Figure 1：用 LLM“显式地”在自然语言里求解\emph{连通性（Connectivity）}任务的例子
  （图内英文为原文保留）。输入是一张无向图和一个问题“节点 7 和节点 9 之间是否有路径？”，
  模型不仅回答“yes”，还写出推理链：节点 7 连到节点 1、节点 1 连到节点 5、节点 5 连到节点 9，
  于是路径为 $[7\to1\to5\to9]$。这种“把思路写出来”正是本文追求的能力。}
  \label{fig:fig1}
\end{figure}

\section{二、摘要中译}

\begin{quote}
大语言模型（LLM）已在众多领域取得令人瞩目的成功，但它们\textbf{理解与求解复杂图问题}的能力却鲜有探索。
为弥补这一空白，我们提出 \textbf{GraphInstruct}——一个新颖的指令微调数据集，旨在让语言模型能够通过
\textbf{显式的推理路径}应对各类图问题。基于 GraphInstruct，我们构建了 \textbf{GraphWiz}：一个开源语言模型，
能够在求解多种图计算问题的同时生成清晰的推理过程。为进一步提升模型的性能与可靠性，我们在图问题求解的
情境中引入了\textbf{直接偏好优化（Direct Preference Optimization, DPO）}框架。改进后的模型 \textbf{GraphWiz-DPO}
在九个不同复杂度的任务上取得了 \textbf{65\%} 的平均准确率，超过了平均准确率为 \textbf{43.8\%} 的 GPT-4。
我们的研究还考察了训练数据量与模型性能之间的关系，强调了随数据量增大而产生的\textbf{过拟合风险}；
并探索了模型在不同任务、不同数据集间的\textbf{可迁移性}，展示了其稳健的零样本泛化能力。GraphWiz 为开发
专精于图推理与图问题求解的 LLM 提供了新的蓝图与有价值的洞见。
\end{quote}

\section{三、三大核心贡献}
\begin{enumerate}[leftmargin=2em]
  \item \textbf{首个大规模图论指令数据集 GraphInstruct}：覆盖 9 类图计算任务，且每条样本都\textbf{带显式推理路径}，
        让训练出的模型既能给答案、又能给过程，大幅提升可解释性。
  \item \textbf{首个图论专精开源模型 GraphWiz}：能解多类型、多规模图问题，性能显著超过当前最强闭源模型 GPT-4。
  \item \textbf{深入的影响因素分析}：细致分析了训练数据量、DPO 中“负样本”采样策略等对性能的影响，
        为后续模型优化提供洞见（尤其揭示了“数据越多不一定越好、存在过拟合阈值”这一现象）。
\end{enumerate}

\section{四、预备知识（小白先看这里）}

\begin{tip}{什么是指令微调（Instruction Tuning）与 SFT}
指令微调 = 用大量“指令（自然语言任务描述）+ 标准答案”的样本去训练 LLM，让它学会\textbf{听懂人话、按要求做事}。
本文用的具体训练方式叫 \textbf{SFT（Supervised Fine-Tuning，监督微调）}：把“图 + 问题”当输入、把“正确的推理路径”
当标准输出，用最普通的“预测下一个词”的语言模型损失来训练。一句话：\textbf{照着标准答案抄}，学会模仿专家的解题写法。
\end{tip}

\begin{tip}{什么是 DPO（直接偏好优化）}
SFT 只教模型“什么是对的”，但不教它“\emph{主动避开}什么是错的”。\textbf{DPO（Direct Preference Optimization）}
就是用来补这一课的。它给模型成对的样本：一个\textbf{更受偏好的}回答 $R_w$（win，正确推理链）和一个\textbf{不受偏好的}
回答 $R_l$（lose，错误推理链），训练目标是让模型\textbf{相对地更喜欢 $R_w$、更不喜欢 $R_l$}。
它是 RLHF（基于人类反馈的强化学习）的一个轻量替代：RLHF 要训练 4 个子模型、又贵又复杂，而 DPO 把它简化成
\emph{一个可以直接优化的损失函数}，不需要单独训练奖励模型。
\end{tip}

\begin{tip}{什么是图论里的“复杂度”与这 9 类任务的难度}
计算机里用\textbf{时间复杂度}衡量一个算法要算多久。本文把 9 类任务按难度分三档：
\textbf{线性复杂度}（$O(|V|+|E|)$，最快，如判连通、找环）、\textbf{多项式复杂度}（稍慢，如最短路、最大流）、
以及 \textbf{NP 完全（NP-Complete）}——目前没有已知的高效算法、规模一大就极难算（如哈密顿路径、子图匹配）。
$|V|$、$|E|$ 分别表示图里节点数、边数。难度越高，越能考验模型是不是“真会推理”。
\end{tip}

\textbf{形式化记号}（看不懂可跳过，用到时再回看）：一张无向或有向图记作 $G=\{V,E\}$，
$V,E$ 分别是节点集与边集。每张图配一个任务相关的提示问题 $Q$。设 $\mathcal{M}$ 为当前的生成式 LLM，
它把每个 $G$-$Q$ 对作为输入，返回一条逐步的文本推理路径 $R$，即 $R=\mathcal{M}(G,Q)$。
这条 $R$ 由多个词或句子组成，在自然语言处理里常被称作\textbf{思维链（Chain-of-Thought, CoT）}。
对某些任务（如最短路），$G$ 里还会附带额外信息，比如边的权重 $w$。

\section{五、数据集 GraphInstruct：9 类任务与“自增强”造数据}

GraphInstruct 是全文的地基。它的两个最大特点是 \textbf{Yes（含训练集）}和 \textbf{Yes（含 CoT 推理路径）}——
这正是它区别于以往 NLGraph、GraphQA 等数据集的地方（后两者都没有训练集、也没有推理路径）。

\subsection{9 类图计算任务（原文 Table 1）}

\begin{table}[H]\centering
\caption{GraphInstruct 的九类任务概览（数据源自原文 Table 1）。$|V|$、$|E|$ 为节点数与边数。}
\begin{tabular}{llccl}
\toprule
\textbf{任务} & \textbf{要做什么} & \textbf{时间复杂度} & \textbf{节点范围} & \textbf{难度} \\
\midrule
Cycle Detection（找环）   & 判断图中是否存在环 & $O(|V|+|E|)$ & [2,100] & 易 \\
Connectivity（连通性）    & 判断两节点间是否有路径 & $O(|V|+|E|)$ & [2,100] & 易 \\
Bipartite Check（二分图） & 判断图是否为二分图 & $O(|V|+|E|)$ & [2,100] & 易 \\
Topological Sort（拓扑排序） & 求有向无环图的拓扑序 & $O(|V|+|E|)$ & [2,50] & 易 \\
Shortest Path（最短路）   & 求两节点间最短路径 & $O(|E|+|V|\log|V|)$ & [2,100] & 中 \\
Maximum Triangle Sum（最大三角形和） & 求权重和最大的连通三元组 & $O(|V|^3)$ & [2,25] & 中 \\
Maximum Flow（最大流）    & 求源点到汇点的最大流量 & $O(|V|^2\sqrt{|E|})$ & [2,50] & 中 \\
Hamilton Path（哈密顿路径） & 判断是否存在恰好经过每点一次的路径 & NP-Complete & [2,50] & 难 \\
Subgraph Matching（子图匹配） & 判断 $G$ 中是否有子图与 $G'$ 同构 & NP-Complete & [2,30] & 难 \\
\bottomrule
\end{tabular}
\end{table}

\subsection{怎么造出来的：六步“自增强（self-augment）”流水线}

直接让人手工写几万条图论推理路径不现实，作者用了一套自动化流程（原文 Figure 2，含 6 步）：
\textbf{① 选任务 $\to$ ② 生成图问题 $\to$ ③ 用 GPT-4 初步生成 $\to$ ④ 训练“数据增强模型” $\to$
⑤ 拒绝采样 $\to$ ⑥ 路径筛选 $\to$ 得到 GraphInstruct}。几个关键点：
\begin{itemize}[leftmargin=2em]
  \item \textbf{随机造图}：用 \textbf{Erd\H{o}s--R\'enyi（ER）随机图模型}（给定节点数 $n$ 和连边概率 $p$）配合
        NetworkX 库批量生成图，并算出标准答案。初始造了 27{,}000 道训练题（9 类各 3{,}000）和 3{,}600 道测试题。
  \item \textbf{先让 GPT-4 起草}：对每道题用 2-shot 思维链提示让 GPT-4 写推理路径，温度设 0.9、重复 3 次以增加多样性，
        \textbf{只保留答对的}，得到初始集 $D_1$（不足 8{,}000 条）。
  \item \textbf{拒绝采样（rejection sampling）补缺}：GPT-4 在某些难任务上答对率很低（如最大流不到 100 条）。
        于是用 $D_1$ 微调一个\textbf{数据增强模型 $\mathcal{M}_1$}（基于 LLaMA 2-13B），让它去\textbf{补答 GPT-4 答错的题}，
        并为同一道题\textbf{生成多条不同的推理路径}。
  \item \textbf{多样性筛选}：用编辑距离、TF-IDF、Jaccard 相似度（基于字符串）和 Sentence-BERT 余弦相似度 +
        K-means 聚类（基于语义）来挑出\textbf{既正确又彼此不同}的推理路径，避免大量雷同。
\end{itemize}

\begin{tip}{为什么“一题配多条推理路径”很重要}
同一道图论题，往往有多种正确的解法走法。让数据集里一道题对应\textbf{多条不同的正确推理链}，能让模型见多识广、
学到多样的解题风格，而不是死记一种套路。最终 GraphInstruct 含 17{,}158 个图-问题对、
共 \textbf{72{,}785} 条推理路径——平均一道题约配 4 条路径。
\end{tip}

\section{六、方法逐节精讲：两阶段训练（SFT $\to$ DPO）}

有了数据，GraphWiz 的训练分\textbf{两个阶段}（见图~\ref{fig:pipe}）：第一阶段用\textbf{混合任务指令微调（SFT）}
让模型学会解题；第二阶段用 \textbf{DPO 对齐}让模型学会“偏好正确链、远离错误链”。骨干模型用开源的
LLaMA 2-7B/13B 和 Mistral-7B。

\begin{figure}[H]\centering
  \includegraphics[width=0.97\linewidth]{cand_p06_k00.pdf}
  \caption{GraphWiz 的两阶段训练流程（原文 Figure 3，图内英文为原文保留）。
  \textbf{左（Phase 1：混合任务指令微调）}：把 9 类任务的“图 + 问题”一起喂给模型做 SFT，
  得到中间模型 $\mathcal{M}_{\mathrm{sft}}$；
  \textbf{右（Phase 2：DPO 对齐）}：给模型成对样本——更受偏好的推理路径 $\mathcal{R}_w$ 与不受偏好的 $\mathcal{R}_l$，
  通过“最大似然偏好”训练，让它更喜欢 $\mathcal{R}_w$，最终得到 GraphWiz。}
  \label{fig:pipe}
\end{figure}

\subsection{第一阶段：混合任务指令微调（SFT）}

\textbf{混合任务（mixed-task）}指的是把 9 类任务\textbf{混在一起同时训练}，而不是一类一类分开训。这样做的好处是：
一个模型通吃所有任务，且任务之间能\textbf{互相促进}（学会找环的能力可能也帮到判连通）。

训练用最普通的\textbf{语言模型损失}：把“图 $G_i$ + 问题 $Q_i$”当输入，让模型逐词预测正确的推理路径 $R_{i,j}$，
最小化预测与真实路径之间的差距：
\begin{equation*}
  \mathcal{L}_{\mathrm{SFT}} = -\sum_{i=1}^{N}\sum_{j=1}^{M}\log \mathbb{P}\!\left(R_{i,j}\mid G_i,Q_i;\theta\right). \tag{2}
\end{equation*}
其中 $N$ 是样本数，$M$ 是 $G_i$-$Q_i$ 对应的推理路径条数，$\theta$ 是模型参数。直观理解：
\textbf{让模型给“正确推理路径”里每个词都打出尽量高的概率}。这一阶段产出的模型记为 $\mathcal{M}(\theta)_{\mathrm{sft}}$。

\subsection{第二阶段：DPO 对齐——本文把 DPO 用到图论上}

SFT 只让模型“模仿对的”，但难任务上模型仍会犯错。DPO 的思路是：\textbf{给它看成对的对错样本，逼它学会区分}。
关键难点是：\textbf{怎么为每道题找到这一对 $R_w$（好）和 $R_l$（坏）？} 作者的做法很巧妙：
\begin{itemize}[leftmargin=2em]
  \item \textbf{$R_w$（偏好/正确链）}：新造 9{,}000 道题，用上一阶段的 $\mathcal{M}_{\mathrm{sft}}$ 对每题推理 20 次，
        选其中\textbf{最长的那条正确路径}作为 $R_w$（长链通常推理更完整详细）。
  \item \textbf{$R_l$（不偏好/错误链）}：故意挑\textbf{“最像正确答案、但其实是错的”那条}（即\emph{难负样本} hard-negative）——
        先用字符串/语义方法找出离 $R_w$ 最近的错误路径，再用\textbf{多数投票（Majority Voting）}选出最常出现的错误链。
        逼模型去分辨这种“似是而非”的陷阱，训练效果更好。
\end{itemize}
最终收集到 6{,}166 个 $G$-$Q$-$R_w$-$R_l$ 四元组。DPO 的训练目标函数为：
\begin{equation*}
  \mathcal{L}_{\mathrm{DPO}}\!\left(\mathcal{M}(\theta);\mathcal{M}(\theta)_{\mathrm{sft}}\right)
  = -\,\mathbb{E}_{(x,R_w,R_l)\sim D}\,
  \log\sigma\!\left(\beta\log\frac{\mathcal{M}(\theta)(R_w\mid x)}{\mathcal{M}(\theta)(R_l\mid x)}
  -\beta\log\frac{\mathcal{M}(\theta)_{\mathrm{sft}}(R_w\mid x)}{\mathcal{M}(\theta)_{\mathrm{sft}}(R_l\mid x)}\right). \tag{3}
\end{equation*}
其中 $x$ 是“图 $G$ 与问题 $Q$ 的拼接”，$\sigma$ 是 sigmoid 函数，$\beta$ 是一个超参数。

\begin{tip}{把 DPO 公式 (3) 拆开讲明白}
公式里有两个分数比，对照着看就懂：
\begin{itemize}[leftmargin=1.5em]
  \item $\dfrac{\mathcal{M}(\theta)(R_w\mid x)}{\mathcal{M}(\theta)(R_l\mid x)}$：\textbf{正在训练的模型}有多偏爱“好链”相对“坏链”。
  \item $\dfrac{\mathcal{M}(\theta)_{\mathrm{sft}}(R_w\mid x)}{\mathcal{M}(\theta)_{\mathrm{sft}}(R_l\mid x)}$：\textbf{冻结的 SFT 模型}的偏好（当作一个“参照基准”）。
\end{itemize}
两者相减、乘以 $\beta$、过 sigmoid，再取负号当损失——\textbf{效果就是：在不偏离原 SFT 模型太远的前提下，
推着新模型相对地更偏好 $R_w$、更冷落 $R_l$}。$\beta$ 控制“偏好对比”的强弱：太大太激进、太小没效果，需要调。
这正是 DPO 优雅之处：\textbf{不用训练奖励模型，一个损失函数就把“偏好对齐”这件事干了}。
\end{tip}

\section{七、实验与结果表解读}

\textbf{对比对象分两大类}：①\textbf{上下文学习（In-Context Learning）}，模型参数不变，靠提示词答题——
代表是闭源的 GPT-3.5、GPT-4（用零样本和 2-shot 思维链）；②\textbf{用简单标签训练（Naive-SFT）}，
只给模型“yes/no”直答、不给推理过程；此外还训了 GCN/GIN/GAT 三种图神经网络（GNN）作对比。

\subsection{主结果：GraphWiz 全面超过 GPT-4（原文 Table 4 关键行重排）}

\begin{table}[H]\centering
\caption{主结果对比（准确率 \%，数值源自原文 Table 4；越高越好，每列加粗为该组最优）。
“Average”为九任务平均。}
\begin{tabular}{llccccc}
\toprule
\textbf{类别} & \textbf{方法} & \textbf{找环} & \textbf{连通} & \textbf{拓扑排序} & \textbf{哈密顿} & \textbf{Average} \\
\midrule
\multirow{3}{*}{闭源模型}
 & GPT-4 (zero-shot) & 38.75 & 17.00 & 5.00 & 59.25 & 27.67 \\
 & GPT-3.5 (2-shot)  & 51.25 & 43.75 & 4.50 & 54.25 & 32.97 \\
 & GPT-4 (2-shot)    & 52.50 & 62.75 & 25.25 & 75.75 & 43.81 \\
\midrule
\multirow{3}{*}{Mistral-7B}
 & Naive SFT      & 73.75 & 83.50 & 1.00 & 31.75 & 46.56 \\
 & GraphWiz       & 92.00 & 89.50 & 19.00 & 26.50 & 53.75 \\
 & GraphWiz-DPO   & 85.50 & 79.50 & \textbf{85.25} & 62.75 & 58.22 \\
\midrule
\multirow{3}{*}{LLaMA 2-7B}
 & Naive SFT      & 73.75 & 83.50 & 4.00 & 69.00 & 44.81 \\
 & GraphWiz       & 91.50 & 87.00 & 18.00 & 52.25 & 55.08 \\
 & \textbf{GraphWiz-DPO} & 89.00 & 82.50 & 46.75 & \textbf{81.50} & \textbf{65.00} \\
\midrule
\multirow{3}{*}{LLaMA 2-13B}
 & Naive SFT      & 73.75 & 83.75 & 0.50 & 59.75 & 44.69 \\
 & GraphWiz       & \textbf{94.75} & 87.00 & 28.00 & 59.00 & 57.39 \\
 & GraphWiz-DPO   & 87.50 & \textbf{88.50} & 72.75 & 46.50 & 63.89 \\
\bottomrule
\end{tabular}
\end{table}

\textbf{怎么读这张表}：
\begin{itemize}[leftmargin=2em]
  \item \textbf{核心数字}：最强的 \textbf{GraphWiz-DPO（LLaMA 2-7B）平均 65.00\%}，
        远超 \textbf{GPT-4 (2-shot) 的 43.81\%}（高出约 21 个百分点）。这就是摘要里“65\% vs 43.8\%”的来历。
        一个 7B 的开源小模型，平均表现把 GPT-4 甩在身后。
  \item \textbf{DPO 确实有用}：在三种骨干上，GraphWiz-DPO 的平均分都比纯 GraphWiz（SFT）更高
        （如 LLaMA 2-7B 从 55.08 升到 65.00）。在\textbf{拓扑排序}这类难任务上提升最明显
        （Mistral 从 19.00 飙到 85.25）。
  \item \textbf{但 DPO 不是“万灵药”}：它在个别任务上反而掉分（如 Mistral 找环从 92.00 降到 85.50），
        说明 DPO 总体提升推理、但对某些任务仍需进一步调参。
  \item \textbf{“写过程”比“只给答案”强}：GraphWiz 远超 Naive-SFT，证明\textbf{让模型显式推理}本身就能提升正确率，
        而不只是为了可解释性。
\end{itemize}

\subsection{其他重要发现}
\begin{itemize}[leftmargin=2em]
  \item \textbf{数据量并非越多越好（RQ2，原文 Table 5）}：随着“图:推理路径”比例 G:R 从 1:1 加到 1:5，
        平均分总体上升（LLaMA 2-7B：46.56 $\to$ 55.08）；但\textbf{个别任务会过拟合}——如 Mistral 的最大三角形和
        从 1:1 的 47.00\% \emph{下降}到 1:5 的 38.75\%。提示：堆数据存在阈值，需配合验证集监控。
  \item \textbf{很能“举一反三”（RQ3）}：作者训了一个只学两类 NP 难任务（哈密顿、子图匹配）的变体 GraphWiz-High，
        发现它在\emph{从没训练过}的简单/中等任务上也能零样本迁移、明显超过原始骨干，体现跨任务泛化能力。
  \item \textbf{图越大越难（RQ4，原文 Figure 6）}：GraphWiz 和 GPT-4 都随图变大而掉分，但 GraphWiz 在大多数规模上
        稳定领先；不过当节点数\textbf{超过 60} 后两者都几乎失效——因为它们主要靠“枚举”推理，规模一大枚举量指数爆炸。
  \item \textbf{案例对比（原文 Figure 8）}：在一张 89 节点的找环题上，GPT-4 直接给出一个很短的错误答案“没有环”，
        而 GraphWiz 正确找出一条具体的环路并写出完整路径，展现了对大规模复杂图的空间推理与记忆能力。
\end{itemize}

\section{八、小白知识点卡片（速查）}
\begin{itemize}[leftmargin=2em]
  \item \textbf{图计算问题}：靠图算法求解的题，如连通性、最短路、找环、拓扑排序、最大流、哈密顿路径等。
  \item \textbf{显式推理路径 / 思维链（CoT）}：模型把“怎么想出来的”一步步写成文字，而不仅给最终答案。
  \item \textbf{指令微调 / SFT}：用“指令 + 标准答案”训练模型听话做事；SFT 是其最常见的监督训练形式。
  \item \textbf{DPO（直接偏好优化）}：给成对的“好/坏”回答，用一个损失函数直接教模型偏好好的、远离坏的；RLHF 的轻量替代。
  \item \textbf{RLHF / PPO}：基于人类反馈的强化学习；PPO 是其常用算法，但要 4 个子模型、成本高。
  \item \textbf{NP 完全（NP-Complete）}：一类目前无已知高效算法、规模一大就极难的问题。
  \item \textbf{ER 随机图}：用“节点数 $n$ + 连边概率 $p$”随机生成图的经典模型。
  \item \textbf{拒绝采样（rejection sampling）}：反复采样、只保留满足条件（这里是“答对”）的结果来扩充数据。
  \item \textbf{难负样本（hard-negative）}：长得最像正确答案、却是错的样本；用它训练能逼模型学得更细。
  \item \textbf{过拟合（overfitting）}：模型死记训练数据的细节，反而在新题上变差。
\end{itemize}

\section{九、一句话总结与谱系定位}

\textbf{一句话}：GraphWiz 先用“GPT-4 起草 + 小模型自增强”造出带显式推理链的图论指令数据集 GraphInstruct，
再\textbf{两阶段训练}——先 SFT 学会解 9 类图论题、写出推理过程，后用 DPO 让模型\textbf{偏好正确链、远离错误链}，
最终用一个 7B 开源模型把平均准确率做到 65\%，反超 GPT-4 的 43.8\%。

\textbf{在“图 $\times$ 后训练”谱系里的位置}：它是\textbf{“指令微调 + 偏好对齐（SFT+DPO）”}路线在\emph{图计算推理}
方向的代表作，核心卖点是\textbf{显式推理 + 开源 + 通吃多任务}。与本系列其他工作对照看更清楚：
G-Retriever 走“检索 + 软提示（GraphRAG）”、LLaGA 走“投影器对齐”、G1 走“强化学习”，
而 GraphWiz（与同期的 InstructGraph）走的是\textbf{“造高质量推理数据 + SFT + DPO”}这条路——
你可以对照着体会\emph{它们各自把“图的推理能力”塞进 LLM 的方式有何不同}。

\textbf{动手复现}：官方代码与数据见 \url{https://github.com/nuochenpku/Graph-Reasoning-LLM}，
项目主页 \url{https://graph-wiz.github.io}。其中 GraphInstruct 数据集本身就是入门“LLM 图推理”很好的练手材料。

\end{document}
